% Adapted from CVQC-paper/figures/compiled-test-messages.tex, retained when
% its inclusion was removed in paper commit 017f09b (2026-11-02).
\begin{frame}{Holler--Natarajan RSP \href{https://arxiv.org/abs/2609.38060}{[HN26]}}
  \centering
  \begin{tikzpicture}[
      >={Latex[length=2.3mm]},
      party/.style={font=\large\bfseries,text=CVQCInk,inner sep=0pt},
      lifeline/.style={densely dashed,draw=CVQCMuted!45,thick},
      message/.style={->,very thick,draw=CVQCTeal},
      message label/.style={above=3pt,fill=white,inner sep=2pt,
        font=\normalsize,text=CVQCInk},
      preparation/.style={rounded corners=2mm,fill=CVQCPale,
        inner xsep=0.25cm,inner ysep=0.18cm,minimum width=3.6cm,
        minimum height=0.80cm,align=center,font=\small,text=CVQCInk}
    ]
    \path[use as bounding box] (-6.6,-3.15) rectangle (6.6,3.05);

    \node[party] at (-4.3,2.65) {Verifier};
    \node[party] at (4.3,2.65) {Prover};
    \draw[lifeline] (-4.3,2.30) -- (-4.3,-0.65);
    \draw[lifeline] (4.3,2.30) -- (4.3,-0.65);

    \draw[message] (-4.3,1.65) -- node[message label]
      {Encrypted question} (4.3,1.65);
    \draw[message] (4.3,0.75) -- node[message label]
      {Encrypted answer} (-4.3,0.75);

    \node[preparation] at (-4.3,-0.15)
      {Compute state description};
    \node[preparation] at (4.3,-0.15)
      {Keep quantum state};

    \only<2>{
      \draw[lifeline] (-4.3,-0.65) -- (-4.3,-2.65);
      \draw[lifeline] (4.3,-0.65) -- (4.3,-2.65);
      \draw[message] (-4.3,-1.70) -- node[message label]
        {Test question} (4.3,-1.70);
      \draw[message] (4.3,-2.50) -- node[message label]
        {Answer} (-4.3,-2.50);
      \node[font=\small,text=CVQCInk,inner sep=0pt] at (-4.3,-2.95)
        {Check};
    }
  \end{tikzpicture}

  \note<1>{
    \footnotesize
    Show the common encrypted exchange first. A preparation request W
    specifies the desired single-qubit bases, and the first question is
    a equals Question of W. Generate fresh encryption keys and send the
    encrypted question a-tilde together with the public evaluation key,
    suppressed in this drawing. The prover returns the classical encrypted
    answer alpha and retains its quantum system. The verifier decrypts
    alpha to v and privately records W,v. Honest preparation leaves the
    corresponding eigenstates on the prover's side. The diagram does not
    assert that a malicious prover's memory already has this ideal form.
    The first-response circuit is the same for preparation and testing.
    The imported construction uses classical-client quantum fully
    homomorphic encryption, with the subexponential quantum-security
    assumption used in our paper. Keep those implementation details off
    the slide. Actual calls prepare nine qubits, one computation qubit
    and eight tags; defer that organization and the conjugation issue.
  }
  \note<2>{
    \footnotesize
    A test continues with two more classical messages: b in the clear
    and the answer gamma. The verifier checks them together with a and
    the decrypted first answer v, using the imported self-test's score
    predicate. The verifier samples the test question pair privately;
    a preparation instead embeds W in the first question. The prover
    does not receive the clear second question before committing to
    its encrypted first answer. Preparation ends after the first two
    messages; only testing uses all four. The test measurements consume
    the prepared state, so it is not retained as an MBQC input on a test
    call. Explain this distinction verbally as the second overlay appears.
    Later slides will explain how many
    preparations and tests yield inputs good enough for our protocol;
    do not claim a negligible-error or composable preparation guarantee
    directly from one invocation. Source: the paper's RSP subprotocol
    interface and its former compiled-test-messages figure, removed
    from the text on November 2 but retained as a figure source.
  }
\end{frame}
